Several validity threats remain.

\paragraph{External validity.}
The study uses public Alibaba traces rather than a live platform, so the conclusions are strongest for trace-driven comparison and weaker for direct deployment claims.

\paragraph{Construct validity.}
The service cohort narrows the construct-validity gap but is not a full service census. It is bounded by a metadata prefilter requiring at least 60 containers and 650{,}000 seconds of span, followed by usage-profile and post-aggregation quality thresholds. The benchmark should therefore be read as a broad rule-based service cohort, not the full \texttt{app\_du} population.

\paragraph{Internal validity.}
\texttt{UA-MSTCN-Lite} is a lightweight trainable surrogate, not the final deep UA-MSTCN architecture. The reported numbers therefore apply to the verified surrogate and policy route, not to an unimplemented deeper model.

\paragraph{Conclusion validity.}
Forecasting gains do not automatically imply better capacity control, so forecasting and policy metrics are kept separate. The service-descriptor splits and taxonomy are descriptive benchmark analyses, not causal identification for live services.
